\subsection{A mixed-unit measure and the exact lifting obstruction}
\label{subsec:cm-nonabelian-gate}

Keep the notation of Subsection~\ref{subsec:cm-elliptic-radical}, and put
\[
 N_{\mathrm{ell}}=\ker\!\left(A\longrightarrow
                  \operatorname{Gal}(L_{\mathrm{ell}}/L_\infty)\right),
 \qquad \overline A=A/N_{\mathrm{ell}},\qquad
 \Lambda_2=\mathbb Z_p[[\Gamma]].
\]
The following construction uses only the mixed torsion points whose
two coordinates are units.  The passage from its moments to the
full radical is justified by an explicit scalar formula, including
the two boundary families of weights.

\begin{lemma}\label{lem:cm-mixed-measure-embedding}
Fix an integer \(c_0>1\) prime to \(6p\), put
\(\mathfrak a=(c_0)\), and set
\[
 u_n=\theta_{\mathfrak a}(\omega_{1,n}+\omega_{2,n}).
\]
There is a continuous injective \(G\)-equivariant homomorphism
\[
 \overline\Phi_{c_0}:\overline A
            \lhook\joinrel\longrightarrow \mathbb Z_p[[G]](1).
\]
Here the target action is
\(\gamma\cdot\mu=\chi_{\mathrm{cyc}}(\gamma)[\gamma]\mu\);
translation is defined by
\([\gamma]\mu(B)=\mu(\gamma^{-1}B)\).
In particular \(\overline A\) is torsion-free over \(\Lambda_2\).
\end{lemma}

\begin{proof}
Let \(G_n=\operatorname{Gal}(L_n/K)\).  The exact norm relation
of \cite[Proposition 2.5, p.~8]{IshiiSoule} gives
\(N_{L_{n+1}/L_n}(u_{n+1})=u_n\).  Its root-of-unity factors are
one because no nonidentity unit of \(O_K\) is congruent to one
modulo \(p\), for \(p\geq5\).
For \(h\in H\), define a bounded measure by the finite-level
Kummer values
\begin{equation}\label{eq:cm-mixed-measure}
 \Phi_{c_0}(h)(\sigma\operatorname{Gal}(L_\infty/L_n))
                      =k_{\sigma(u_n)}(h)\qquad(\sigma\in G_n).
\end{equation}
The norm identity says exactly that the value of a cylinder is
the sum of the values of its refinements.  Thus these values give
an element of \(\mathbb Z_p[[G]]\).  Additivity of Kummer characters
makes \(\Phi_{c_0}\) a continuous homomorphism on \(A\).
For a lift of \(\gamma\in G\), the defining root equation gives
\[
 k_u(\gamma h\gamma^{-1})
   =\chi_{\mathrm{cyc}}(\gamma)k_{\gamma^{-1}u}(h).
\]
This proves the stated action and its independence of the lift,
since inner conjugation acts trivially on \(A\).

We identify its kernel using the moment formula for all weights.  By the proof of
Corollary~\ref{cor:cm-radical-joint-kernel}, a positive integral
power of each \(u_n\) belongs to \(C(L_n)\).  Consequently
\(N_{\mathrm{ell}}\subseteq\ker\Phi_{c_0}\).
Conversely, suppose \(\Phi_{c_0}(h)=0\).
For every \(\boldsymbol m=(m_1,m_2)\in I\), the exact moment
formula of \cite[Lemma 4.2, p.~17]{IshiiSoule} gives
\[
 \kappa_{\boldsymbol m,\mathfrak a}(h)
   =\int_G\chi^{\boldsymbol m-\boldsymbol1}(\sigma)\,
                                 d\Phi_{c_0}(h)(\sigma)=0.
\]
In this source \(\kappa_{\boldsymbol m,\mathfrak a}\) is defined
from this same mixed norm system for every \(\boldsymbol m\in I\).
When one of \(m_1,m_2\) equals one, the proof of
\cite[Proposition 4.3, pp.~17--18]{IshiiSoule} takes the norm
to the corresponding one-prime ray tower and includes its Euler
factor.  Its conclusion still reads
\begin{equation}\label{eq:cm-mixed-all-weight-scalar}
 w_K\,\frac{1-p^{r+s}}{(1-A_{\boldsymbol m})(1-B_{\boldsymbol m})}
       \kappa_{\boldsymbol m,\mathfrak a}
    =(c_0^2-c_0^{-r-s})\kappa_{\boldsymbol m},
 \quad r=m_1-1,\quad s=m_2-1,
\end{equation}
where
\[
 A_{\boldsymbol m}=i^{(r,s)}(\pi),\qquad
 B_{\boldsymbol m}=i^{(r,s)}(\bar\pi),\qquad
 A_{\boldsymbol m}B_{\boldsymbol m}=p^{r+s}.
\]
These algebraic numbers have complex absolute value
\(p^{(r+s)/2}>1\), because \(r,s\geq0\) and \(r+s>0\).
Both denominators and the numerator on the left are therefore
nonzero, including \(r=0\) and \(s=0\).
The right scalar is also nonzero.  The unit ambiguity in the
Artin action disappears since \(r-s\equiv0\pmod{w_K}\).
Thus \(\kappa_{\boldsymbol m}(h)=0\) for every
\(\boldsymbol m\in I\).
Corollary~\ref{cor:cm-radical-joint-kernel} now gives
\(h\bmod[H,H]\in N_{\mathrm{ell}}\).
This proves the exact equality
\(\ker\Phi_{c_0}=N_{\mathrm{ell}}\).

Finally \(G/\Gamma\) is finite of order prime to \(p\), so the
untwisted \(\mathbb Z_p[[G]]\) is finite free over \(\Lambda_2\).
The cyclotomic twist changes scalars by the continuous ring
automorphism \(\gamma\mapsto\chi_{\mathrm{cyc}}(\gamma)\gamma\)
of \(\Lambda_2\).  Its target remains torsion-free, and the
displayed injection proves the assertion.
\end{proof}

\begin{lemma}\label{lem:cm-exact-eigenlift-obstruction}
Let \(\gamma\neq1\) belong to \(\Gamma\), let
\(\lambda\in\mathbb Z_p^\times\), and let \(t\) be a lift of
\(\gamma\) to \(\operatorname{Gal}(\Omega/K)\).
If an element \(h\in H\) satisfies
\[
                         tht^{-1}=h^\lambda ,
\]
then every elliptic Soul\'e character vanishes on \(h\).
If \(\mathcal X_2^+\) is finite, then \(h\in[H,H]\).
\end{lemma}

\begin{proof}
Write \(a\) for the image of \(h\) in \(\overline A\).
The asserted group equation gives the exact module equation
\[
                              (\gamma-\lambda)a=0.
\]
The element \(\gamma-\lambda\) is nonzero in
\(\Lambda_2=\mathbb Z_p[[\mathbb Z_p^2]]\), an integral domain.
Lemma~\ref{lem:cm-mixed-measure-embedding} therefore gives \(a=0\).
The joint-kernel description implies the first assertion.
Under finite \(\mathcal X_2^+\),
Theorem~\ref{thm:cm-finite-plus-abelian} gives
\(N_{\mathrm{ell}}=0\), hence \(h\) has zero image in \(A\).
\end{proof}

This gives an obstruction to lifting homogeneous
Galois eigenclasses by group elements satisfying the corresponding
eigen-equations.  Rational eigenspace decompositions in a graded
quotient do not provide such equations in \(H\).

\begin{proposition}\label{prop:cm-no-normalized-cyclic}
Assume that \(\mathcal X_2^+\) is finite.  Put
\(G_1=\operatorname{Gal}(\Omega/L_1)\), so that
\[
             1\longrightarrow H\longrightarrow G_1
               \longrightarrow\Gamma\longrightarrow1 .
\]
If \(t\in G_1\) has nonidentity image in \(\Gamma\), then the
equation
\[
                 tht^{-1}=h^\lambda,\qquad
                    h\in H,\quad\lambda\in\mathbb Z_p^\times
\]
has only the solution \(h=1\).
Consequently, for every nontrivial closed procyclic subgroup
\(C\subset H\), the image of \(N_{G_1}(C)\) in \(\Gamma\) is
trivial.
\end{proposition}

\begin{proof}
The finite-plus theorem gives that \(H\) is free pro-\(p\) and
that \(A=H^{\mathrm{ab}}\) embeds equivariantly in
\[
 P=\mathbb Z_p[[G]](1).
\]
As a \(\Lambda_2\)-module, \(P\) is a finite direct sum of
copies of the regular module, with its cyclotomic twist.

We first justify taking completed tensor powers of this
injection.  Every compact torsion-free \(\mathbb Z_p\)-module
is topologically a product of copies of \(\mathbb Z_p\).
Indeed, its Pontryagin dual is a divisible \(p\)-primary group,
and hence a direct sum of copies of
\(\mathbb Q_p/\mathbb Z_p\).  For such a product \(M\), and any
compact \(\mathbb Z_p\)-module \(V\), the definition by finite
continuous quotients gives
\[
 \left(\prod_{i\in I}\mathbb Z_p\right)
       \widehat\otimes_{\mathbb Z_p}V
             \simeq\prod_{i\in I}V.
\]
This functor preserves injections.  Applying this assertion
one tensor factor at a time gives, for every \(d\geq1\), an
injection
\[
 A^{\widehat\otimes d}\lhook\joinrel\longrightarrow
 P^{\widehat\otimes d}
     \simeq\mathbb Z_p[[G^d]](d).
\]

The action of \(\Gamma\) on the last module is diagonal
translation with the scalar twist
\(\chi_{\mathrm{cyc}}^d\).  The continuous group isomorphism
\[
 \Gamma\times\Gamma^{d-1}\longrightarrow\Gamma^d,\qquad
 (g,h_2,\ldots,h_d)\longmapsto
                (g,gh_2,\ldots,gh_d)
\]
identifies diagonal translation with translation in the first
factor.  After choosing the finitely many tame cosets, its
completed group algebra is therefore a finite direct sum of
power-series modules
\[
             \Lambda_2[[T_1,\ldots,T_{2d-2}]].
\]
It is torsion-free over \(\Lambda_2\).  The scalar twist changes
the coefficient action by the ring automorphism
\(\gamma\mapsto\chi_{\mathrm{cyc}}(\gamma)^d\gamma\), which
preserves that property.  Thus every completed tensor power
of \(A\) is \(\Lambda_2\)-torsion-free.

Write \(\gamma_dH\) for the closed lower central series.
For a free pro-\(p\) group, the canonical map
\[
 \gamma_dH/\gamma_{d+1}H
        \longrightarrow A^{\widehat\otimes d},
 \qquad
 [h_1,\ldots,h_d]\longmapsto
          [\bar h_1,\ldots,\bar h_d]_{\mathrm{tensor}},
\]
identifies its source with the completed degree-\(d\) free
Lie submodule.  This assertion is integral: for a finite free
basis, the Hall commutators give the free
\(\mathbb Z_p\)-basis, and their Magnus expansions have the
corresponding leading Lie tensors.  For a general free
pro-\(p\) basis converging to one, projection onto finite
subsets and reduction modulo \(p^n\) give the completed
assertion.  These projections also prove injectivity, since a
nonzero coefficient is seen in a finite set of basis
coordinates and modulo some \(p^n\).  The construction is
functorial in automorphisms.  In particular, the conjugation
action on this graded piece is the action induced by that on
\(A\), and so it factors through \(\Gamma\).
It follows that every
\(\gamma_dH/\gamma_{d+1}H\) is
\(\Lambda_2\)-torsion-free.

Every pro-\(p\) group is separated for its closed lower central
series: any nonidentity element survives in a finite \(p\)-group,
which is nilpotent.  Therefore, if \(h\neq1\), it has a nonzero
image \(a\) in some
\(\gamma_dH/\gamma_{d+1}H\).  If \(t\) has image
\(\gamma\neq1\), the proposed equation gives
\[
                         (\gamma-\lambda)a=0 .
\]
The power series \(\gamma-\lambda\) is nonzero in
\(\Lambda_2\), contradicting the torsion-freeness just proved.
Thus \(h=1\).

Finally \(H\) is torsion-free, so a nontrivial closed
procyclic \(C\) is isomorphic to \(\mathbb Z_p\).
Conjugation by a normalizer element sends a topological
generator \(h\) to \(h^\lambda\) for a unit \(\lambda\).
The first assertion forces the image of that element in
\(\Gamma\) to be one.  This proves the precise normalizer
claim.
\end{proof}

The proposition gives a further constraint on a proposed
arithmetic use of the geometric recognition theorem.  A
nontrivial procyclic subgroup selected inside \(H\), including
one selected inside the arithmetic outer kernel, cannot
acquire an open normalizer image in
\(\operatorname{Gal}(L_\infty/L_1)\): its normalizer image
there is exactly the identity subgroup.  This conclusion
concerns the displayed arithmetic extension and follows from
its actual conjugation action on every lower-central quotient.

\paragraph{The remaining correction equation.}
Let \(\widetilde F^\bullet H\) be the weighted free filtration
attached to a proposed free set of arithmetic lifts, with its
filtered outer map to \(F^\bullet\rho(H)\).
A homogeneous derivation relation, such as the relation in
bidegree \((12,2)\) proved above, produces a word
\(r\in\widetilde F^{14}H\) satisfying
\[
 \rho(r)\in F^{15}\rho(H).
\]
To cancel that outer image without changing its initial
free class, one must solve the stronger, explicit equation
\[
 \rho(r)=\rho(s)\quad\hbox{for some }s\in\widetilde F^{15}H.
\]
Its solution gives \(rs^{-1}\in\ker\rho\).
The first inclusion alone does not supply \(s\):
it concerns the target filtration, whereas the displayed equation
asks for membership in the image of a specified source subgroup.
Rational surjectivity of successive graded maps gives no
uniform integral bound on the denominators in successive
corrections.  Clearing a growing denominator by successive
\(p\)-powers can erase the proposed initial class.
The argument in \cite[proof of Theorem 2.14, p.~27]{Ishii} uses graded
freeness precisely to identify these two filtrations.
The source-group freeness established here does not itself
solve the displayed correction equation.
